% ----------------------------------------------------------
% Pacotes e Configurações
% ----------------------------------------------------------

% ---
% Pacotes básicos
% ---
\usepackage{enumitem}
\usepackage{lmodern}      % Usa a fonte Latin Modern
\usepackage[T1]{fontenc}    % Selecao de codigos de fonte.
\usepackage[utf8]{inputenc}   % Codificacao do documento (conversão automática dos acentos)
\usepackage{lastpage}     % Usado pela Ficha catalográfica
\usepackage{indentfirst}    % Indenta o primeiro parágrafo de cada seção.
\usepackage{color}        % Controle das cores
\usepackage{graphicx}     % Inclusão de gráficos
\usepackage{microtype}    % para melhorias de justificação
\usepackage{float}
\usepackage{caption}
\usepackage{subcaption}
\usepackage{booktabs}
% ---

% ---
% Pacotes adicionais, usados apenas no âmbito do Modelo Canônico do abnteX2
% ---
\usepackage[geometry]{ifsym}
\usepackage{multicol}
\usepackage{pdfpages}
\usepackage[printonlyused]{acronym} % pacote para produzir acrônimos - utilize [printonlyused] para gerar a lista somente com os que foram utilizados

% --- Definição das siglas para o pacote acronym ---
\acrodef{AP}{Access Point}
\acrodef{API}{Application Programming Interface}
\acrodef{BYOD}{Bring Your Own Device}
\acrodef{CCQ}{Client Connection Quality}
\acrodef{COSI}{Colegiado do Curso de Sistemas de Informação}
\acrodef{CSMA/CA}[CSMA/CA]{Carrier Sense Multiple Access with Collision Avoidance}
\acrodef{DECSI}{Departamento de Computação e Sistemas}
\acrodef{DHCP}{Dynamic Host Configuration Protocol}
\acrodef{HTTP}{Hypertext Transfer Protocol}
\acrodef{ICEA}{Instituto de Ciências Exatas e Aplicadas}
\acrodef{IEEE}{Institute of Electrical and Electronics Engineers}
\acrodef{JM}{João Monlevade}
\acrodef{JSON}{JavaScript Object Notation}
\acrodef{LGPD}{Lei Geral de Proteção de Dados}
\acrodef{MAC}{Medium Access Control}
\acrodef{MCS}{Modulation and Coding Scheme}
\acrodef{MIB}{Management Information Base}
\acrodef{MIMO}{Multiple-Input Multiple-Output}
\acrodef{MU-MIMO}[MU-MIMO]{Multi-User Multiple-Input Multiple-Output}
\acrodef{OFDMA}{Orthogonal Frequency-Division Multiple Access}
\acrodef{OID}{Object Identifier}
\acrodef{PHY}{Physical Layer}
\acrodef{QoS}{Quality of Service}
\acrodef{REST}{Representational State Transfer}
\acrodef{RF}{Radiofrequência}
\acrodef{RSSI}{Received Signal Strength Indicator}
\acrodef{SI}{Sistemas de Informação}
\acrodef{SNMP}{Simple Network Management Protocol}
\acrodef{SNR}{Signal-to-Noise Ratio}
\acrodef{SSID}{Service Set Identifier}
\acrodef{UFOP}{Universidade Federal de Ouro Preto}
\acrodef{UTC}{Coordinated Universal Time}
\acrodef{VLAN}{Virtual Local Area Network}
\acrodef{VRP}{Vehicle Routing Problem}
\acrodef{WLAN}{Wireless Local Area Network}
\acrodef{WPA2}{Wi-Fi Protected Access 2}

% ---

% ---
% Pacotes de citações
% ---
\usepackage[brazilian,hyperpageref]{backref}   % Paginas com as citações na bibl
\usepackage[alf]{abntex2cite} % Citações padrão ABNT

% FBO: Review
\usepackage{color}
\newcommand{\review}[1]{{\textbf{\color{red}{#1}}}}

\newenvironment{reviewblock}
{\bfseries\color{red}}
{\normalfont\color{black}}

% ---
% CONFIGURAÇÕES DE PACOTES
% ---

% ---
% Configurações do pacote backref
% Usado sem a opção hyperpageref de backref
\renewcommand{\backrefpagesname}{Citado na(s) página(s):~}
% Texto padrão antes do número das páginas
\renewcommand{\backref}{}
% Define os textos da citação
\renewcommand*{\backrefalt}[4]{
  \ifcase #1 %
    Nenhuma citação no texto.%
  \or
    Citado na página #2.%
  \else
    Citado #1 vezes nas páginas #2.%
  \fi}%
% ---

% ---

% Tamanho da linha de assinatura segundo o tamanho do maior nome
% Sugestão: 8cm
\setlength{\ABNTEXsignwidth}{8cm}

% ---
% Configurações de aparência do PDF final

% alterando o aspecto da cor azul
\definecolor{blue}{RGB}{41,5,195}

% informações do PDF
\makeatletter
\hypersetup{
      %pagebackref=true,
    pdftitle={\@title},
    pdfauthor={\@author},
      pdfsubject={\imprimirpreambulo},
      pdfcreator={LaTeX with abnTeX2},
    pdfkeywords={abnt}{latex}{abntex}{abntex2}{trabalho acadêmico},
    colorlinks=true,          % false: boxed links; true: colored links
      linkcolor=blue,           % color of internal links
      citecolor=blue,           % color of links to bibliography
      filecolor=magenta,          % color of file links
    urlcolor=blue,
    bookmarksdepth=4
}
\makeatother
% ---

\usepackage{listings}
\usepackage{xcolor}

% Configurações para exibição de código (listings)
\definecolor{codegreen}{rgb}{0,0.6,0}
\definecolor{codegray}{rgb}{0.5,0.5,0.5}
\definecolor{codepurple}{rgb}{0.58,0,0.82}
\definecolor{backcolour}{rgb}{0.95,0.95,0.92}

\lstdefinestyle{mystyle}{
    backgroundcolor=\color{backcolour},   
    commentstyle=\color{codegreen},
    keywordstyle=\color{magenta},
    numberstyle=\tiny\color{codegray},
    stringstyle=\color{codepurple},
    basicstyle=\ttfamily\footnotesize,
    breakatwhitespace=false,         
    breaklines=true,                 
    captionpos=b,                    
    keepspaces=true,                 
    numbers=left,                    
    numbersep=5pt,                  
    showspaces=false,                
    showstringspaces=false,
    showtabs=false,                  
    tabsize=2
}
\lstset{style=mystyle}
% ---

% ---
% Espaçamentos entre linhas e parágrafos
% ---

% O tamanho do parágrafo é dado por:
\setlength{\parindent}{1.3cm}

% Controle do espaçamento entre um parágrafo e outro:
\setlength{\parskip}{0.2cm}  % tente também \onelineskip

% ---
% compila o indice
% ---
\makeindex
% ---
